\chapter{Die Stärkeschätzung: ELO}\label{kap:elo}

Das ELO-System stammt aus dem Schach. Jede Mannschaft trägt eine Zahl, die
ihre Spielstärke schätzt. Nach jedem Spiel wird verglichen, was das Modell
erwartet hat und was passiert ist. Die Differenz, verstärkt durch die Höhe
des Ergebnisses, verschiebt die Zahlen beider Mannschaften um denselben
Betrag in entgegengesetzte Richtungen.

In der Bundesliga liegen die Werte grob zwischen 1\,500 und 2\,100, in den
Regionalligen etwa zwischen 700 und 1\,200. Die Skala hat keinen absoluten
Nullpunkt. Bedeutung haben nur Unterschiede, und das auch nur innerhalb einer
Wechselgemeinschaft, also einer Gruppe von Ligen, zwischen denen Mannschaften
auf- und absteigen (Abschnitt~\ref{sec:wechselgemeinschaft}).

Die Rechnung steht in einer einzigen Funktion,
\code{calculate\_elo\_change}, in
\datei{league-simulator-rust/src/elo/mod.rs:5--32}. Sie übernimmt die Logik
der früheren C++-Implementierung (\datei{SpielNichtSimulieren.cpp}).

\section{Erwartungswert}

Es seien $E_h$ und $E_a$ die ELO-Werte der Heim- und der Gastmannschaft
vor dem Spiel und $H = 40$ der Heimvorteil. Der Code bildet zuerst die
\emph{invertierte} Differenz und begrenzt sie auf $\pm 400$:
\begin{equation}
  \Delta_{\text{inv}} = E_a - E_h - H, \qquad
  \tilde\Delta = \clamp_{[-400,\,400]}\bigl(\Delta_{\text{inv}}\bigr)
              = \min\bigl(\max(\Delta_{\text{inv}}, -400),\, 400\bigr).
\end{equation}
Der Erwartungswert der Heimmannschaft ist die klassische logistische
ELO-Formel mit Basis 10 und Skala 400:
\begin{equation}\label{eq:elo-erwartung}
  p_h = \frac{1}{1 + 10^{\tilde\Delta/400}}.
\end{equation}

\begin{itemize}
  \item Der Heimvorteil wird auf die ELO der Heimmannschaft aufgeschlagen.
        Bei gleich starken Mannschaften erwartet das Modell für die
        Heimmannschaft $p_h = 1/(1+10^{-0{,}1}) \approx 0{,}557$.
  \item Wegen der Begrenzung gilt immer
        $p_h \in [\tfrac{1}{11},\, \tfrac{10}{11}] \approx [0{,}091;\ 0{,}909]$.
        Auch ein krasser Favorit kann bei einem Sieg also noch Punkte gewinnen.
  \item $p_h$ ist ein \emph{erwarteter Score}, bei dem ein Remis als halber
        Sieg zählt. Es ist \emph{nicht} die Siegwahrscheinlichkeit des
        Tormodells (Abschnitt~\ref{sec:skalen}).
\end{itemize}

\section{Tatsächlicher Score und Tordifferenz-Multiplikator}

Das Ergebnis wird auf einen Score abgebildet. Ein Sieg zählt 1, ein Remis
0{,}5 und eine Niederlage 0
(\datei{elo/mod.rs:16--17}):
\begin{equation}
  S = \frac{\sgn(g_h - g_a) + 1}{2} \in \{0,\ \tfrac12,\ 1\}.
\end{equation}
Hohe Siege zählen mehr, aber gedämpft. Der Multiplikator ist die Wurzel der
Tordifferenz, mindestens 1 (\datei{elo/mod.rs:19--20}):
\begin{equation}
  m = \sqrt{\max\bigl(\lvert g_h - g_a\rvert,\ 1\bigr)}.
\end{equation}
Remis und Siege mit einem Tor Unterschied haben $m = 1$. Ein 4:0 hat
$m = 2$ und verschiebt damit doppelt so viel wie ein 1:0.

\section{Das Update}

Mit dem Gewichtungsfaktor $k = 20$ (\code{mod\_factor}) gilt
(\datei{elo/mod.rs:23--27}):
\begin{equation}\label{eq:elo-update}
  \boxed{\;\delta = k \cdot m \cdot (S - p_h), \qquad
  E_h' = E_h + \delta, \qquad E_a' = E_a - \delta\;}
\end{equation}

\begin{description}[font=\sffamily\bfseries\small]
  \item[Nullsumme] Was die eine Mannschaft gewinnt, verliert die andere. In
        einer geschlossenen Liga bleiben Summe und Mittelwert der ELO über die
        ganze Saison konstant. Der Rust-Test
        \code{elo\_is\_conserved\_by\_the\_walk} prüft das, und die Eichung
        neuer Ligen (Abschnitt~\ref{sec:eichung}) nutzt es.
  \item[Keine Rundung] ELO ist durchgehend eine Gleitkommazahl. Die
        Stammdatei speichert volle Genauigkeit, zum Beispiel
        \code{2057.25650898366} für den FC Bayern. Gerundet wird nur bei der
        Anzeige.
  \item[Reihenfolge] Die Updates werden in chronologischer Reihenfolge der
        Anstoßzeiten ausgeführt, ausgehend vom ELO-Wert zum Saisonstart
        (\code{InitialELO} in \datei{RCode/TeamList_2026.csv}). Diese
        Kette heißt im Code \emph{ELO-Walk}.
  \item[Gewertete Spiele] Spiele, die am grünen Tisch gewertet wurden
        (Status AWD oder WO), zählen für die Tabelle, bewegen die ELO aber
        nicht (Issue~\#157). Ein 3:0 am grünen Tisch sagt nichts über die
        Spielstärke aus.
\end{description}

\section{Rechenbeispiel}\label{sec:elo-beispiel}

Das Eröffnungsspiel der Saison 2026/27: FC Bayern ($E_h = 2057{,}26$) gegen
VfB Stuttgart ($E_a = 1805{,}06$).
\begin{align*}
  \Delta_{\text{inv}} &= 1805{,}06 - 2057{,}26 - 40 = -292{,}19
    \quad\text{(innerhalb $\pm 400$)},\\
  p_h &= 1/(1 + 10^{-292{,}19/400}) = 0{,}8432.
\end{align*}

\begin{center}
\begin{tabular}{@{}lccr@{}}
\toprule
Ergebnis & $S$ & $m$ & ELO-Änderung Bayern $\delta$ \\
\midrule
2:1 & 1   & 1          & $20 \cdot 1 \cdot (1 - 0{,}8432) = +3{,}14$ \\
4:0 & 1   & 2          & $+6{,}27$ \\
1:1 & 0{,}5 & 1        & $-6{,}86$ \\
0:2 & 0   & $\sqrt2$   & $-23{,}85$ \\
\bottomrule
\end{tabular}
\end{center}

Stuttgart erhält jeweils den Gegenwert mit umgekehrtem Vorzeichen. Ein Remis
kostet den Favoriten mehr, als ein knapper Sieg ihm einbringt: Das Modell
hatte einen Sieg weitgehend vorausgesetzt.

Eine zweite Gegenprobe steht in ADR~0002. Ein 3:0 bei 1\,500 gegen 1\,480
ergibt $\delta = 20\sqrt3\,(1 - p_h) = +14{,}359$.

\begin{merke}[Wie die ELO-Werte auf die Seite kommen]
Gespeichert wird nur die ELO zum Saisonstart. Den aktuellen Stand und jeden
Zwischenstand rechnet der Rust-Server bei jedem Rendern neu aus allen
gespielten Spielen nach. Das geschieht deterministisch über den Endpunkt
\code{POST /league-details} (ADR~0002). Ligatabelle, Rückblick, Ausblick und
ELO-Verlauf beruhen so alle auf demselben Walk und können sich nicht
widersprechen.
\end{merke}
